\documentclass[../main]{subfiles}
\begin{document}
\chapter{Costos Fijos}
\img{images/01/costo.jpg}{Costos}
\section{Costos Fijos}
\info{¿Qué son los costos fijos?}{
  Los costos fijos son aquellos cuyo pago es asumido por la empresa de
  manera constante, independientemente de su participación dentro del
  proceso productivo. A estos costos se les conoce como fijos porque no
  varían ante los cambios de la producción de bienes y servicios. Es
  decir, son aquellos que, sin importar cuánto se produzca, siempre
  deberán ser abonados. Por ejemplo: el alquiler de una oficina o un
  local, los sueldos administrativos, los servicios de telefonía e
internet, el pago de seguros, etc.}
\subsection{¿Cómo determinarlos?}
Los costos fijos pueden determinarse acorde a las decisiones de la
administración, de tal forma que esta puede prescindir de un servicio
o, por el contrario, adquirir uno, o preservar los ya existentes segñun
sean las necesidades de la organización.
\subsection{Importancia}
La importancia de llevar un control de estos costos radica en su gran
incidencia en la determinación de los precios finales de los bienes o
servicios que pueda proporcionar una empresa. Además, deben ser
incluidos dentro de los presupuestos por ser costos que se presentan
de forma constante y permanente.

\subsection{Características}

A continuación, se relacionan las principales características de
los costos fijos:
\begin{enumerate}
  \item No varían por la presencia de cambios en la producción.
  \item Pueden determinarse en función de las decisiones administrativas.
  \item A pesar de estar ligados a la capacidad de producción que
    posee una organización, estos no sufren variaciones si se llegan
    a alterar los niveles de producción.
  \item Cumplen un rol fundamental para el cálculo del punto de equilibrio.
    Se debe llevar un registro detallado de estos costos para poder
    tener un control oportuno sobre ellos.
\end{enumerate}
\begin{center}
  \begin{tikzpicture}
    \draw (0,0)--(0,5)[-stealth] node [left] {C};
    \draw (0,0)--(5,0)[-stealth] node [below] {Q};
    \draw (0,2)--(5,2) node [below] {CF: Costos Fijos};
  \end{tikzpicture}
\end{center}
\subsection{Tipos de costos fijos}
\subsubsection{Discrecionales}
Aquellos que pueden llegar a reducirse sin que la producción de la
organización sea afectada. Estos costos se originan con la finalidad
de ser empleados en actividades o áreas determinadas y las
organizaciones suelen ser más flexibles sobre estos. Los costos
discrecionales generalmente se originan en planes tácticos de
corto alcance, generalmente para un marco de tiempo de un año o más.
Por ejemplo: campañas de publicidad, el alquiler de ciertos
 bienes, investigaciones, entre otros.
\subsubsection{Comprometidos}
Aquellos costos que no son susceptibles de sufrir modificaciones
porque el proceso productivo de la empresa puede ser afectado. Por
ende, los costos fijos comprometidos corresponden a costos necesarios
y de carácter obligatorio ya que, si estos son interrumpidos, pueden
presentarse inconvenientes severos en la rentabilidad de la
organización. Por el contrario, la interrupción de los procesos de
producción por algñun factor externo a estos gastos no alterará
los costos fijos comprometidos, debido a la naturaleza que presentan.

Los costos comprometidos se originan en planificaciones estratégicas
que generalmente abarcan un marco temporal de varios años y que en
gran medida fijan el nivel y el tipo de operaciones de una empresa.
Estos costos se deben controlar de manera cuidadosa con un adecuado
análisis previo, para optimizar de la manera más eficiente
los elementos que integran los costos fijos comprometidos.

Por ejemplo: salarios de empleados, seguros, impuestos sobre raíz de
propiedad, entre otros.

\subsection{Ejemplos de costos fijos}
\begin{itemize}
  \item  El pago del alquiler de un local, bodega u oficina.
  \item Pagos por seguros en caso de siniestros.
  \item Depreciación de maquinaria y equipo.
  \item Costos de investigación o desarrollo.
  \item Pagos por servicios no medibles de telefonía, internet, agua, etc.
  \item Amortización de costos de patentes, marcas y licencias.
  \item Los permisos municipales y licencias.
  \item Los salarios de los empleados administrativos.
\end{itemize}

\newpage
\section{El Punto de Equilibrio}
\info{Concepto}{El \textbf{Punto de Equilibrio} es el nivel de producción o ventas en el cual los ingresos totales son iguales a los costos totales (tanto fijos como variables). En este punto, la empresa no obtiene ganancias ni sufre pérdidas.}

Para calcular el punto de equilibrio en unidades ($Q$), se utiliza la siguiente fórmula:

\begin{equation}
  Q = \frac{CF}{P - CV}
\end{equation}

Donde:
\begin{itemize}
  \item \textbf{$Q$}: Cantidad de equilibrio (unidades a producir/vender).
  \item \textbf{$CF$}: Costos Fijos Totales.
  \item \textbf{$P$}: Precio de Venta Unitario.
  \item \textbf{$CV$}: Costo Variable Unitario.
\end{itemize}

La diferencia $(P - CV)$ se conoce como el \textbf{Margen de Contribución unitario}, que es la cantidad de dinero de cada venta que contribuye a cubrir los costos fijos.

\ejemplo{Ejemplo Rápido}{Un taller de mecanizado tiene costos fijos de \$1000 mensuales. El precio de un servicio de fresado es de \$50 y el costo variable por pieza es de \$30. El punto de equilibrio es: $Q = 1000 / (50 - 30) = 50$ piezas. El taller debe realizar 50 servicios para no perder dinero.}

\actividad{Taller de Análisis de Costos: La Startup Electromecánica}

\textbf{Escenario}: Ustedes han decidido abrir un pequeño taller de servicios industriales. Cuentan con los siguientes activos y gastos mensuales:
\begin{itemize}
  \item Alquiler del local: \$500
  \item Seguro contra incendios: \$50
  \item Sueldo de un administrativo (part-time): \$300
  \item Amortización de un Torno CNC: \$150
  \item Publicidad en redes sociales: \$100
\end{itemize}

\textbf{Consignas:}
\begin{enumerate}
  \item \textbf{Categorización}: Clasifiquen cada uno de los gastos anteriores como Costos Fijos \textit{Comprometidos} o \textit{Discrecionales}.
  \item \textbf{Cálculo de CF}: Determinen el Costo Fijo Total mensual del taller.
  \item \textbf{Análisis de Equilibrio}: Si el precio de venta de un servicio promedio es de \$120 y el costo variable (materiales y energía) es de \$70, calculen el Punto de Equilibrio en unidades.
  \item \textbf{Análisis Crítico}: Si el alquiler del local aumenta un 20\%, \textbf{¿cómo afecta esto al punto de equilibrio?} Realicen el nuevo cálculo y expliquen la consecuencia para el negocio.
\end{enumerate}
\end{document}